% !TEX TS-program = pdflatex
% !TeX program = pdflatex
% !TEX encoding = UTF-8
% !TEX spellcheck = fr

\documentclass[xcolor=table]{beamer}


%\usepackage{fullpage}
%\usepackage[left=2.8cm,right=2.2cm,top=2 cm,bottom=2 cm]{geometry}
\setbeamersize{text margin left=10pt,text margin right=10pt}


\usepackage{amsmath,amssymb} 
\usepackage{fontenc}
\usepackage{textcomp}
\usepackage[utf8]{inputenc}
\usepackage{CJKutf8}
\usepackage[french]{babel}
\usepackage{arabtex,acolor}
\usepackage{txfonts}
\usepackage{tipa}
\usepackage[]{graphicx}
\usepackage{multirow}
\usepackage{hyperref}
\usepackage{colortbl}
\usepackage{tabularray}
\usepackage{listingsutf8}
\usepackage{wrapfig}
\usepackage{multicol}
\usepackage[export]{adjustbox} %for images in table, also for frame
\usepackage[many]{tcolorbox}
\usepackage{wasysym}
\usepackage[french,lined]{algorithm2e}
\usepackage{alltt} %verbatim with commands
\usepackage{longtable}
\usepackage{tabu}


\definecolor{lightblue}{HTML}{D0D2FF}
\definecolor{lightyellow}{HTML}{FFFFAA}
\definecolor{darkblue}{HTML}{0000BB}
\definecolor{olivegreen}{HTML}{006600}
\definecolor{darkgreen}{HTML}{008B45} %009B55
\definecolor{violet}{HTML}{6600CC}
%\definecolor{deeppink}{HTML}{FF1493}
\definecolor{orangey}{HTML}{FFBB00}

\usetheme{Warsaw} % Antibes Boadilla Warsaw

\beamertemplatenavigationsymbolsempty


\let\oldcite\cite
\renewcommand{\cite}[1]{{\bfseries\color{orangey}\oldcite{#1}}}

\tcbuselibrary{listings}

\newcommand{\kurl}[1]{{\scriptsize\bfseries\color{orangey}\url{#1}}}


%\renewcommand{\baselinestretch}{1.5}

\def\supit#1{\raisebox{0.8ex}{\small\it #1}\hspace{0.05em}}

\AtBeginSection{%
	\begin{frame}
		\sectionpage
	\end{frame}
}

\newcommand{\rottext}[2]{%
	\rotatebox{90}{%
	\begin{minipage}{#1}%
		\raggedleft#2%
	\end{minipage}%
	}%
}


\institute{ %
Laboratoire de la Communication dans les Systèmes Informatiques (LCSI)

\vspace{6pt}
École  nationale Supérieure d'Informatique (ESI, ex. INI), Alger, Algérie
}
\author[ \textbf{\footnotesize\insertframenumber/\inserttotalframenumber} \hspace*{1cm} ESI - ARIES Abdelkrime (Master 2022/2023)] %
{ARIES Abdelkrime}
%\titlegraphic{\includegraphics[height=1cm]{../img/esi-logo.png}%\hspace*{4.75cm}~
	
	
	\date{Master 2022/2023} %\today

\titlegraphic{%
	\includegraphics[height=1cm]{../img/esi-logo.png}%
	\hspace{2cm}
	\includegraphics[height=1cm]{../img/lcsi-logo.png}%
}


%\setbeamertemplate{headline}{}



\newcommand{\keyword}[1]{\textcolor{red}{\bfseries\itshape #1}}
\newcommand{\expword}[1]{\textcolor{olivegreen}{#1}}
\newcommand{\optword}[1]{\textcolor{violet}{\bfseries #1}}

\makeatletter
\newcommand\mysphere{%
	\parbox[t]{10pt}{\raisebox{0.2pt}{\beamer@usesphere{item projected}{bigsphere}}}}
\makeatother

%\let\oldtabular\tabular
%\let\endoldtabular\endtabular
%\renewenvironment{tabular}{\rowcolors{2}{white}{lightblue}\oldtabular\rowcolor{blue}}{\endoldtabular}


\NoAutoSpacing %french autospacing after ":"

\def\graphpath{}

\newcommand{\changegraphpath}[1]{\def\graphpath{#1}}


\newcommand{\vgraphpage}[2][.84\textheight]{%
%	\begin{center}%
		\includegraphics[height=#1]{\graphpath #2}%
%	\end{center}%
}

\newcommand{\hgraphpage}[2][\textwidth]{%
%	\begin{center}%
		\includegraphics[width=#1]{\graphpath #2}%
%	\end{center}%
}

\newcommand{\graphpage}[2][]{%
	\includegraphics[#1]{\graphpath #2}%
}

\bibliographystyle{apalike}

\newcommand{\insertbibliography}[2]{
	\appendix
	\section*{Bibliographie}
	\nocite{#2}
%	\makeatletter % to change template
%	\setbeamertemplate{headline}[default] % not mandatory, but I though it was better to set it blank
%	\def\beamer@entrycode{\vspace*{-\headheight}} % here is the part we are interested in :)
%	\makeatother
	\begin{multicols*}{2}[\frametitle{\insertsection} \usebeamertemplate*{frametitle}]%\usebeamertemplate*{frametitle}\frametitle{Références}
		\tiny
		\bibliography{#1}
	\end{multicols*}
}

\definecolor{my-grey}{RGB}{233, 233, 233}

\newcommand{\insertlicence}{
	\begin{frame}[plain]
	\frametitle{Licence : CC-BY 4.0}
%	\framesubtitle{Licence: CC-BY-NC 4.0}

	\begin{tcolorbox}[colback=cyan,
		colframe=cyan,  
		arc=0pt,outer arc=0pt,
		valign=top, 
		halign=center,
		width=\textwidth]
		
		\includegraphics[width=.5cm]{../img/licence/cc_icon_white_x2.png}
		\includegraphics[width=.5cm]{../img/licence/attribution_icon_white_x2.png}
		
		\color{white}
		\bfseries Attribution 4.0 International (CC BY 4.0) \\
		\tiny \url{https://creativecommons.org/licenses/by/4.0/deed.fr}
		
	\end{tcolorbox}\vspace{-.5cm}
	\begin{tcolorbox}[colback=my-grey,
		colframe=my-grey,  
		center, arc=0pt,outer arc=0pt,
		valign=top, 
		halign=left,
		width=\textwidth]
		
		\tiny
		
		\begin{center}
			\bfseries\Large
			Vous êtes autorisé à :
		\end{center}
		
		\begin{minipage}{0.83\textwidth}
			\begin{itemize}
				\item[] \textbf{Partager} — copier, distribuer et communiquer le matériel par tous moyens et sous tous formats
				\item[] \textbf{Adapter} — remixer, transformer et créer à partir du matériel
				pour toute utilisation, y compris commerciale.
			\end{itemize}
		\end{minipage}
		\begin{minipage}{0.15\textwidth}
			\includegraphics[width=\textwidth]{../img/licence/FreeCulturalWorks_seal_x2.jpg}
		\end{minipage}
	
		
		\begin{center}
			\bfseries\Large
			Selon les conditions suivantes :
		\end{center}
		
		\begin{itemize}
			\item[] \textbf{Attribution} — Vous devez créditer l'Œuvre, intégrer un lien vers la licence et indiquer si des modifications ont été effectuées à l'Oeuvre. Vous devez indiquer ces informations par tous les moyens raisonnables, sans toutefois suggérer que l'Offrant vous soutient ou soutient la façon dont vous avez utilisé son Oeuvre. 
			\item[] \textbf{Pas de restrictions complémentaires} — Vous n'êtes pas autorisé à appliquer des conditions légales ou des mesures techniques qui restreindraient légalement autrui à utiliser l'Oeuvre dans les conditions décrites par la licence.
		\end{itemize}
		
	\end{tcolorbox}
	
%	\begin{center}
%		\bfseries Attribution 4.0 International (CC BY 4.0)
%		\url{https://creativecommons.org/licenses/by/4.0/deed.fr}
%	\end{center}

%	\tiny
%
%	Vous êtes autorisé à : 
%	\begin{itemize}
%		\item \textbf{Partager} — copier, distribuer et communiquer le matériel par tous moyens et sous tous formats
%		\item \textbf{Adapter} — remixer, transformer et créer à partir du matériel
%	\end{itemize}
%	
%	Selon les conditions suivantes : 
%	\begin{itemize}
%		\item \textbf{Attribution} — Vous devez créditer l'Œuvre, intégrer un lien vers la licence et indiquer si des modifications ont été effectuées à l'Oeuvre. Vous devez indiquer ces informations par tous les moyens raisonnables, sans toutefois suggérer que l'Offrant vous soutient ou soutient la façon dont vous avez utilisé son Oeuvre.
%		\item \textbf{Pas d'Utilisation Commerciale} — Vous n'êtes pas autorisé à faire un usage commercial de cette Oeuvre, tout ou partie du matériel la composant. 
%		\item \textbf{Pas de restrictions complémentaires} — Vous n'êtes pas autorisé à appliquer des conditions légales ou des mesures techniques qui restreindraient légalement autrui à utiliser l'Oeuvre dans les conditions décrites par la licence.
%	\end{itemize}

	\end{frame}
}

\settowidth{\leftmargini}{\usebeamertemplate{itemize item}}
\addtolength{\leftmargini}{\labelsep}

\AtBeginDocument{
	\newcolumntype{L}[2]{>{\vbox to #2\bgroup\vfill\flushleft}p{#1}<{\egroup}} 
	
	\begin{frame}[plain]
		\maketitle
	\end{frame}

	\insertlicence
}


% needs etoolbox; to break links after -
\appto\UrlBreaks{\do\-}


\makeatletter
\newcommand{\xRightarrow}[2][]{\ext@arrow 0359\Rightarrowfill@{#1}{#2}}
\makeatother


\usefonttheme{structurebold}
%\usefonttheme{professionalfonts}

\title[TALN : 04- Étiquetage morpho-syntaxique]%
{Traitement automatique du langage naturel\\Chapitre 04 : Étiquetage morpho-syntaxique} 

\changegraphpath{../img/pos/}

\begin{document}
	
\begin{frame}
\frametitle{Traitement automatique du langage naturel}
\framesubtitle{Étiquetage morpho-syntaxique : Introduction}

\begin{itemize}
	\item We can can the can
	
\end{itemize}

\end{frame}


\begin{frame}
\frametitle{Traitement automatique du langage naturel}
\framesubtitle{Étiquetage morpho-syntaxique : Plan}

\begin{multicols}{2}
%	\small
\tableofcontents
\end{multicols}
\end{frame}

%===================================================================================
\section{Étiquetage de séquences}
%===================================================================================

\subsection{Description}

\begin{frame}
\frametitle{Étiquetage morpho-syntaxique : Séquences}
\framesubtitle{Description}


\begin{itemize}
	\item \optword{But :} affecter des étiquettes à des éléments d'une séquence. 
	\item \optword{Exemple : } \expword{affecter des catégories grammaticales à des mots dans une phrase}.
\end{itemize}

\begin{block}{Formulation du problème}
	\begin{itemize}
		\item $w = w_1, \ldots, w_n$ : une séquence d'entrée
		\item $t = t_1, \ldots, t_n$ : une séquence des étiquettes
	\end{itemize}
	\begin{center}
		$ \arg\max\limits_t P(t | w)$
	\end{center}
	
	\begin{itemize}
		\item \optword{Modèle génératif} $ \arg\max\limits_t P(t | w) = \arg\max\limits_t P(t) P(w | t) $
		\item \optword{Modèle discriminatif} calculer $\arg\max\limits_t P(t | w)$ directement
	\end{itemize}
\end{block}

\end{frame}

\subsection{Applications}

\begin{frame}
\frametitle{Étiquetage morpho-syntaxique : Séquences}
\framesubtitle{Applications}

%eisenstein
\begin{itemize}
	\item Reconnaissance d'entités nommées
	\begin{itemize}
		\item Trouver les personnes, les organisations, les places, les nombres, etc. dans une phrase
	\end{itemize}
	\item Étiquetage morpho-syntaxique
	\begin{itemize}
		\item Trouver les catégories des mots (nom, verbe, etc.) d'une phrase
	\end{itemize}
	\item Analyse syntaxique de surface (Chunking)
	\begin{itemize}
		\item Trouver les syntagmes d'une phrase sans structure d'arbre
	\end{itemize}
\end{itemize}


\end{frame}

\begin{frame}
\frametitle{Étiquetage morpho-syntaxique : Séquences}
\framesubtitle{Applications (exemple)}

\begin{figure}
	\centering
	\hgraphpage{exp-ner_.pdf}
	\caption{Un exemple de la reconnaissance d'entités nommées [\url{https://corenlp.run/}]}
\end{figure}

\begin{figure}
	\centering
	\hgraphpage{exp-pos_.pdf}
	\caption{Un exemple d'étiquetage morpho-syntaxique [\url{https://corenlp.run/}]}
\end{figure}

\end{frame}

%\subsection{Évaluation}
%
%\begin{frame}
%\frametitle{Étiquetage morpho-syntaxique : Séquences}
%\framesubtitle{Évaluation}
%
%\begin{itemize}
%	\item 
%\end{itemize}
%
%\end{frame}

%===================================================================================
\section{Description de la tâche}
%===================================================================================

\begin{frame}
\frametitle{Étiquetage morpho-syntaxique}
\framesubtitle{Description de la tâche}


\begin{itemize}
	\item \optword{Objectif} : trouver les catégories grammaticales de chaque mot dans une phrase
	\item Définir les catégories grammaticales
	\begin{itemize}
		\item Nature, catégorie grammatical, catégorie lexicale, classe grammaticale, espèce grammaticale, partie du discours
		\item Anglais : Part of speech 
	\end{itemize}
	\item Annoter un corpus
	\begin{itemize}
		\item Annoter un texte manuellement pour le test
		\item Un autre pour l'entraînement si on utilise l'apprentissage automatique ou on veut fixer des hyper-paramètres
	\end{itemize}
\end{itemize}

\end{frame}

\subsection{Classes universelles}

\begin{frame}
\frametitle{Étiquetage morpho-syntaxique : Description de la tâche}
\framesubtitle{Classes universelles}

\url{https://universaldependencies.org/u/pos/}

\rowcolors{2}{lightblue}{lightyellow}
\begin{tabular}{p{.3\textwidth}p{.3\textwidth}p{.3\textwidth}}
	\rowcolor{darkblue}
	\textcolor{white}{Classe ouverte} & \textcolor{white}{Classe fermée} & \textcolor{white}{Autres} \\
	
	\keyword{ADJ} :  adjectif & \keyword{ADP} : adposition & \keyword{PUNCT} : ponctuation \\
	\keyword{ADV} :  adverbe & \keyword{AUX} : auxiliaire & \keyword{SYM} : symbole \\
	\keyword{INTJ} : interjection & \keyword{CCONJ} : conjonction de coordination & \keyword{X} : Autres \\
	\keyword{NOUN} : nom & \keyword{DET} : déterminant &  \\
	\keyword{PROPN} : nom propre & \keyword{NUM} : numérique &  \\
	\keyword{VERB} : verbe & \keyword{PART} : particule &  \\
	 & \keyword{PRON} : pronom &  \\
	 & \keyword{SCONJ} : conjonction de subordination &  \\
	
\end{tabular}

\end{frame}

\subsection{Treebanks}

\begin{frame}[fragile]
\frametitle{Étiquetage morpho-syntaxique : Description de la tâche}
\framesubtitle{Treebanks}


\begin{figure}
	\begin{tcolorbox}
		\small
\begin{alltt}
Battle-tested\keyword{/JJ} Japanese\keyword{/JJ} industrial\keyword{/JJ} managers\keyword{/NNS}
here\keyword{/RB} always\keyword{/RB} buck\keyword{/VBP} up\keyword{/RP} nervous\keyword{/JJ} newcomers\keyword{/NNS}
with\keyword{/IN} the\keyword{/DT} tale\keyword{/NN} of\keyword{/IN} the\keyword{/DT} first\keyword{/JJ} of\keyword{/IN}
their\keyword{/PP\$} countrymen\keyword{/NNS} to\keyword{/TO} visit\keyword{/VB} Mexico\keyword{/NNP} ,\keyword{/,}
a\keyword{/DT} boatload\keyword{/NN} of\keyword{/IN} samurai\keyword{/FW} warriors\keyword{/NNS} blown\keyword{/VBN}
ashore\keyword{/RB} 375\keyword{/CD} years\keyword{/NNS} ago\keyword{/RB} .\keyword{/.}
\end{alltt}
\end{tcolorbox}
\caption{Exemple d'un texte annoté de Penn TreeBank \cite{2003-taylor}}
\end{figure}

\begin{itemize}
	\item Universal TreeBank (\url{https://universaldependencies.org/}) \cite{2012-petrov-al}
	\begin{itemize}
		\item Projet à communauté ouverte
		\item Il vise à fournir des corpus annotés pour plusieurs langues
		\item Il utilise les mêmes classes grammaticales
	\end{itemize}
\end{itemize}

\end{frame}

\subsection{Difficulté et outils}

\begin{frame}
\frametitle{Étiquetage morpho-syntaxique : Description de la tâche}
\framesubtitle{Difficulté et outils : Difficulté}

\begin{itemize}
	\item \optword{Ambigüité}
	\begin{itemize}
		\item Il existe des mots avec plusieurs catégories grammaticales
		\item Ex. \expword{We can can the can}
	\end{itemize}
	\item \optword{Mots inconnus}
	\begin{itemize}
		\item On peut trouver des mots hors vocabulaire (du corpus d'entraînement)
%		\item Ex. \expword{We can can the can}
	\end{itemize}
	\item \optword{Inconsistance des tags} 
	\begin{itemize}
		\item Lorsqu'il y a plusieurs tags, on peut trouver des inconsistances dans l'annotation
		\item Ex. \expword{Dans les corpus ``Brown" et ``WSJ", le mot ``to" a le tag ``TO". Dans le corpus ``Switchboard", il a le tag ``TO" lorsqu'il indique l'infinitif, sinon ``IN" lorsque c'est une préposition.}
	\end{itemize}
\end{itemize}

\end{frame}

\begin{frame}
\frametitle{Étiquetage morpho-syntaxique : Description de la tâche}
\framesubtitle{Difficulté et outils : Outils}

\begin{itemize}
	\item \url{https://nlp.stanford.edu/software/tagger.shtml}
	\item \url{https://github.com/sloria/textblob}
	\item \url{https://www.nltk.org/api/nltk.tag.html}
	\item \url{https://spacy.io/}
	\item \url{https://github.com/clips/pattern}
\end{itemize}

\end{frame}

%===================================================================================
\section{Approches}
%===================================================================================

\begin{frame}
\frametitle{Étiquetage morpho-syntaxique}
\framesubtitle{Approches}

\begin{itemize}
	\item 
\end{itemize}

\end{frame}

\subsection{Modèle de Markov}

\begin{frame}[fragile]
\frametitle{Étiquetage morpho-syntaxique : Approches}
\framesubtitle{Modèle de Markov}

\begin{minipage}{.64\textwidth}
	\begin{itemize}
		\item Hypothèse de Markov
		\item $ P(q_i = a | q_1, \ldots, q_{i-1}) \approx P(q_i = a | q_{i-1}) $
		\item $Q = \{q_1, q_2, \ldots, q_n\}$ : états
		\item $A = \begin{bmatrix}%
		a_{11} & a_{12} & \ldots & a_{1n} \\
		a_{21} & a_{22} & \ldots & a_{2n} \\
		\vdots & \vdots & \ddots & \vdots \\
		a_{n1} & a_{n2} & \ldots & a_{nn} \\
		\end{bmatrix}$ : matrice des probabilités de transition
		\item $\sum_j a_{ij} = 1,\, \forall i$
	\end{itemize}
\end{minipage}
\begin{minipage}{.35\textwidth}
	\begin{figure}
		\hgraphpage{exp-markov_.pdf}
		\caption{Exemple d'un chaîne de Markov pour la prédiction du temps \cite{2019-jurafsky-martin}}
	\end{figure}
\end{minipage}

\begin{itemize}
	\item $\pi = [\pi_1, \pi_2, \ldots, \pi_n ]$ : la distribution initiale des probabilités des états
	\item $\sum_i \pi_i = 1$
	\item Ex. \expword{Calculer la probabilité $P(HOT\, WARM\, COLD\, COLD)$ lorsque $\pi = [0.1, 0.7, 0.2]$}
\end{itemize}

\end{frame}


\begin{frame}[fragile]
\frametitle{Étiquetage morpho-syntaxique : Approches}
\framesubtitle{Modèle de Markov : Modèle de Markov caché (HMM)}

\begin{minipage}{.54\textwidth}
	\begin{itemize}
%		\item Hypothèse de Markov
%		\item $ P(q_i = a | q_1, \ldots, q_{i-1}) \approx P(q_i = a | q_{i-1}) $
		\item $Q = \{q_1, q_2, \ldots, q_n\}$ : états
		\item $A$ : matrice des probabilités de transitions où $\sum_j a_{ij} = 1,\, \forall i$
		\item $O = o_1 o_2 \ldots o_T$ : séquence des évènements observés (mots)
		\item $B = b_i(o_t)$ : probabilités d'observation (\keyword{probabilités d'émission}), chacune représente la probabilité de génération d'une observation $o_t$ dans un état $q_i$
	\end{itemize}
\end{minipage}
\begin{minipage}{.45\textwidth}
	\begin{figure}
		\hgraphpage{exp-hmm_.pdf}
		\caption{Exemple d'une représentation HMM \cite{2019-jurafsky-martin}}
	\end{figure}
\end{minipage}

\begin{itemize}
	\item $\pi = [\pi_1, \pi_2, \ldots, \pi_n ]$ : la distribution initiale des probabilités des états
	\item $\sum_i \pi_i = 1$
\end{itemize}

\end{frame}

\begin{frame}
\frametitle{Étiquetage morpho-syntaxique : Approches}
\framesubtitle{Modèle de Markov : Entraînement}

\begin{itemize}
	\item Les états sont les tags (catégories) $t_i$
	\item Les observation $o_i$ sont les mots $w_i$
\end{itemize}

\begin{block}{Entraînement du HMM pour étiquetage morpho-syntaxique}
	\[
	P(t_i | t_{i-1}) = \frac{C(t_{i-1}, t_i)}{t_{i-1}} 
	\hskip2cm
	P(w_i | t_i) = \frac{C(t_i, w_i)}{t_i}
	\]
	
\end{block}

\end{frame}

\begin{frame}
\frametitle{Étiquetage morpho-syntaxique : Approches}
\framesubtitle{Modèle de Markov : Entraînement (exemple)}

\begin{exampleblock}{Exemple d'un corpus d'entraînement}
	\begin{itemize}
		\item un/DET ordianteur/NOUN peut/VERB vous/PRON aider/VERB
		\item il/PRON veut/VERB vous/PRON aider/VERB
		\item il/PRON veut/VERB un/DET ordinateur/NOUN
		\item il/PRON peut/VERB nager/VERB
	\end{itemize}
\end{exampleblock}

\begin{itemize}
	\item $P(VERB | NOUN) = \frac{C(NOUN,\ VERB)}{C(NOUN)} = \frac{1}{2}$
	\item $P(veut | VERB) = \frac{C(VERB,\ veut)}{C(VERB)} = \frac{2}{7}$
\end{itemize}

\end{frame}

\begin{frame}
\frametitle{Étiquetage morpho-syntaxique : Approches}
\framesubtitle{Modèle de Markov : Étiquetage (décodage des tags)}

\begin{itemize}
	\item Étant donné 
	\begin{itemize}
		\item un modèle de Markov caché $\lambda = (A, B)$
		\item une séquence des observations (mots) : $w = w_1 w_2 \ldots w_n$
	\end{itemize}
	\item Estimer la séquence des étiquettes $\hat{t} = \hat{t}_1 \hat{t}_2 \ldots \hat{t}_n$
\end{itemize}

\begin{block}{Décodage des étiquètes en utilisant HMM}
	\[
	\hat{t} = \arg\max\limits_t P(t | w) = \arg\max\limits_t \frac{P(w|t) P(t)}{P(w)} = \arg\max\limits_t P(w|t) P(t)%\text{ tel que } t = t_1 t_2 \ldots t_n
	\]
	
	\[ 
	P(w | t) \approx \prod\limits_{i=1}^n P(w_i|t_i) 
	\hskip2cm
	P(t) \approx \prod\limits_{i=1}^n P(t_i|t_{i-1}) 
	\]
	
	\[
	\hat{t} = \arg\max\limits_t \prod\limits_{i=1}^n P(w_i|t_i) P(t_i|t_{i-1})
	\]
\end{block}

\end{frame}

\begin{frame}
\frametitle{Étiquetage morpho-syntaxique : Approches}
\framesubtitle{Modèle de Markov : Viterbi}
\vspace{-6pt}
\begin{block}{Viterbi }
	\scriptsize
	\begin{algorithm}[H]
		\KwData{$w = w_1 \ldots w_T$, HMM $\lambda = (A, B)$ avec $N$ états}
		\KwResult{$meilleur\_chemin$, $prob\_chemin$}
		
		Créer une matrice $viterbi[N, T]$\;
		
		\Pour{état $ s = 1 \ldots N$}{
			$viterbi[s, 1] = \pi_s * b_s(w_1);\, backpointer[s, 1] = 0$ \;
		}
		
		\Pour{$ t = 1 \ldots T$}{
			\Pour{état $ s = 1 \ldots N$}{
				$viterbi[s, t] = \max\limits_{s'=1}^N viterbi[s', t-1] * a_{s',s} * b_s(w_t)$\;
				$backpointer[s, t] = \arg\max\limits_{s'=1}^N viterbi[s', t-1] * a_{s',s} * b_s(w_t)$\;
			}
		}
	
		$prob\_chemin = \max\limits_{s=1}^N viterbi[s, T];\, pointeur\_chemin = \arg\max\limits_{s=1}^N viterbi[s, T]$\;
		
		$meilleur\_chemin$ est le chemin qui commence par $pointeur\_chemin$ et qui suit $backpointer$
		
		\Return $meilleur\_chemin$, $prob\_chemin$\;
		
	\end{algorithm}
\end{block}

\end{frame}


\begin{frame}
\frametitle{Étiquetage morpho-syntaxique : Approches}
\framesubtitle{Modèle de Markov : Exercice}

\begin{exampleblock}{Exemple d'un corpus d'entraînement}
	\begin{itemize}
		\item un/DET ordianteur/NOUN peut/VERB vous/PRON aider/VERB
		\item il/PRON veut/VERB vous/PRON aider/VERB
		\item il/PRON veut/VERB un/DET ordinateur/NOUN
		\item il/PRON peut/VERB nager/VERB
	\end{itemize}
\end{exampleblock}

\begin{itemize}
	\item $G = \{DET, NOUN, VERB, PRON\}$
	\item Trouver les tags de cette phrase : \expword{il peut aider}
	\begin{itemize}
		\item Entraîner le modèle (seulement sur les mots et les tags nécessaires)
		\item Appliquer Viterbi
	\end{itemize}
\end{itemize}

\end{frame}

\subsection{Maximum entropy}

\begin{frame}
\frametitle{Étiquetage morpho-syntaxique : Approches}
\framesubtitle{Maximum entropy : Maximum entropy Markov model (MEMM)}

\begin{itemize}
	\item On note $x_{n}^{m} = x_n x_{n+1} \ldots x_{m-1} x_m$
	\item Étant donné 
	\begin{itemize}
		\item une séquence des observations (mots) : $w = w_1 w_2 \ldots w_n$
		\item un ensemble des caractéristiques $f$ sur les séquences (comme : mot en majuscule)
	\end{itemize}
	\item Estimer la séquence des étiquettes $\hat{t} = \hat{t}_1 \hat{t}_2 \ldots \hat{t}_n$
\end{itemize}

\begin{block}{Décodage des étiquètes en utilisant MEMM}
	\[
	\hat{t} = \arg\max\limits_t P(t | w) = \arg\max\limits_t \prod\limits_{i}  P(t_i | w_i, t_{i-1})
	\]
	
	\[
	\hat{t} = \arg\max\limits_t \prod\limits_{i}  
	\frac{exp\left(\sum_j \theta_j f_j(t_i, w_{i-l}^{i+l}, t_{i-k}^{i-1})\right)}%
	{\sum_{t' \in tags} exp\left(\sum_j \theta_j f_j(t'_i, w_{i-l}^{i+l}, t_{i-k}^{i-1})\right)}
	\]

\end{block}

\end{frame}

\begin{frame}
\frametitle{Étiquetage morpho-syntaxique : Approches}
\framesubtitle{Maximum entropy : Les caractéristiques}

\begin{itemize}
	\item $w_i$ contient un préfixe à partir d'une liste ($len(prefixe) \le 4$) 
	\item $w_i$ contient un suffixe à partir d'une liste ($len(suffixe) \le 4$) 
	\item $w_i$ contient un nombre 
	\item $w_i$ contient une lettre en majuscule
	\item $w_i$ contient un trait d'union 
	\item $w_i$ est complètement en majuscule
	\item La forme du mot $w_i$ (Ex. \expword{X.X.X}) 
	\item $w_i$ est en majuscule avec un trait et un nombre (Ex. \expword{CFC-12}) 
	\item $w_i$ est en majuscule suivi des mots Co., Inc., etc. dans 3 mots maximum
\end{itemize}

\end{frame}

\subsection{Modèle neuronal}

\begin{frame}
\frametitle{Étiquetage morpho-syntaxique : Approches}
\framesubtitle{Modèle neuronal}

\begin{itemize}
	\item 
\end{itemize}

\end{frame}

\insertbibliography{TALN04}{*}

\end{document}

